\documentclass[10pt,twocolumn]{article}

% ============================================================
% PACKAGES
% ============================================================
\usepackage[utf8]{inputenc}
\usepackage[T1]{fontenc}
\usepackage[bahasa]{babel}
\usepackage{geometry}
\usepackage{graphicx}
\usepackage{amsmath}
\usepackage{amssymb}
\usepackage{booktabs}
\usepackage{array}
\usepackage{tabularx}
\usepackage{multirow}
\usepackage{hyperref}
\usepackage{url}
\usepackage{cite}
\usepackage{enumitem}
\usepackage{fancyhdr}
\usepackage{titlesec}
\usepackage{abstract}
\usepackage{microtype}
\usepackage{setspace}
\usepackage{parskip}
\usepackage{caption}
\usepackage{float}
\usepackage{xcolor}

% ============================================================
% PAGE GEOMETRY
% ============================================================
\geometry{
  a4paper,
  top=2.5cm,
  bottom=2.5cm,
  left=2cm,
  right=2cm,
  columnsep=0.8cm
}

% ============================================================
% HYPERREF SETUP
% ============================================================
\hypersetup{
  colorlinks=true,
  linkcolor=blue,
  citecolor=blue,
  urlcolor=blue
}

% ============================================================
% SECTION FORMATTING
% ============================================================
\titleformat{\section}
  {\normalfont\bfseries\large}{\thesection.}{0.5em}{}
\titleformat{\subsection}
  {\normalfont\bfseries\normalsize}{\thesubsection}{0.5em}{}
\titleformat{\subsubsection}
  {\normalfont\bfseries\small}{\thesubsubsection}{0.5em}{}

% ============================================================
% HEADER / FOOTER
% ============================================================
\pagestyle{fancy}
\fancyhf{}
\renewcommand{\headrulewidth}{0.4pt}
\fancyhead[L]{\small\textit{Prototype Monitoring dan Alert System DRC Pemerintah Kota Bogor}}
\fancyhead[R]{\small\thepage}

% ============================================================
% ABSTRACT FORMATTING
% ============================================================
\newenvironment{abstractfull}{%
  \begin{minipage}{\textwidth}
  \small
}{%
  \end{minipage}
}

% ============================================================
% BEGIN DOCUMENT
% ============================================================
\begin{document}

% ---- TITLE (one-column span) --------------------------------
\twocolumn[{%
\begin{@twocolumnfalse}

  \begin{center}
    {\LARGE\bfseries Prototype Monitoring and Alert System for Disaster Recovery Center (DRC) of Bogor City Government}\\[6pt]

    {\large\itshape Prototype Sistem Monitoring dan Alert Disaster Recovery Center (DRC)\\
    Pemerintah Kota Bogor}\\[10pt]

    {\normalsize
      Tsabitha Naylasafa Aurora$^{1}$,
      Faiz Naufal Huda$^{2}$,
      Daffa Kautsar Falih$^{3}$,\\
      Muhammad Naufal Fathurrahman$^{4}$,
      Dr.\ Shelvie Nidya Neyman$^{5}$
    }\\[4pt]

    {\small
      $^{1,2,3,4}$Program Studi Ilmu Komputer, Sekolah Sains Data, Matematika, dan Informatika,\\
      Institut Pertanian Bogor, Indonesia\\[2pt]
      $^{5}$Departemen Ilmu Komputer, Institut Pertanian Bogor, Indonesia\\[2pt]
      \textit{corresponding author:}
      \href{mailto:tsabitha06nayla@apps.ipb.ac.id}{tsabitha06nayla@apps.ipb.ac.id}
    }\\[6pt]

    {\footnotesize
      Copyright \copyright\ 2026 Tsabitha Naylasafa Aurora, et al.
      This is an open-access article distributed under the Creative Commons Attribution License,
      which permits unrestricted use, distribution, and reproduction in any medium,
      provided the original work is properly cited.
    }
  \end{center}

  \vspace{6pt}
  \hrule
  \vspace{6pt}

  % ---- ABSTRACT (English) ----
  \noindent\textbf{Abstract}\\[4pt]
  \noindent\small
  Digital services of Bogor City Government increasingly depend on reliable information technology
  infrastructure, as mandated by Presidential Regulation No.\ 95 of 2018 on Electronic-Based
  Government Systems (SPBE). However, no standardized real-time monitoring mechanism currently exists
  for government-owned Disaster Recovery Centers (DRC) at the regional level. This research designs
  and implements a prototype real-time monitoring and alerting system for a DRC using an open-source
  technology stack: Prometheus, Grafana 12.4.1, Alertmanager, and Node Exporter 1.10.2. The system
  monitors seven critical infrastructure parameters across four virtual machines running Ubuntu Server
  24.04 LTS in a simulated DC-DRC environment, and delivers automatic notifications via Telegram Bot
  API. An inhibit rules mechanism suppresses redundant resource alerts when a NodeDown event is
  active, preventing alert storms. All components run as systemd services with a 15-second scrape
  interval. Testing was conducted through load simulation using \texttt{iperf3} and \texttt{stress-ng}
  in a VMware Workstation environment. The results show that the system successfully detects and
  delivers alerts within $[\mathit{X}]$ seconds of the incident occurring, with threshold detection
  accuracy of $[\mathit{Y}]$\%.\\[4pt]

  \noindent\textbf{Keywords:} alert, disaster recovery center, Grafana, inhibit rules, monitoring,
  Prometheus, systemd, virtual machine.

  \vspace{8pt}
  \hrule
  \vspace{6pt}

  % ---- ABSTRAK (Indonesian) ----
  \noindent\textbf{Abstrak}\\[4pt]
  \noindent\small
  Layanan digital Pemerintah Kota Bogor semakin bergantung pada infrastruktur teknologi informasi
  yang handal, sebagaimana diamanatkan oleh Peraturan Presiden Nomor 95 Tahun 2018 tentang Sistem
  Pemerintahan Berbasis Elektronik (SPBE). Namun, belum terdapat mekanisme pemantauan
  \textit{real-time} yang terstandarisasi untuk \textit{Disaster Recovery Center} (DRC) milik
  pemerintah daerah. Penelitian ini merancang dan mengimplementasikan \textit{prototype} sistem
  \textit{monitoring} dan \textit{alerting real-time} untuk DRC menggunakan \textit{stack} teknologi
  \textit{open-source}: Prometheus, Grafana 12.4.1, Alertmanager, dan Node Exporter 1.10.2. Sistem
  memantau tujuh parameter kritis pada empat \textit{virtual machine} Ubuntu Server 24.04 LTS yang
  mensimulasikan lingkungan DC-DRC, dan mengirimkan notifikasi otomatis melalui Telegram Bot API.
  Mekanisme \textit{inhibit rules} menekan \textit{alert} sumber daya yang redundan saat kondisi
  NodeDown aktif, mencegah \textit{alert storm}. Seluruh komponen berjalan sebagai layanan
  \textit{systemd} dengan \textit{scrape interval} 15 detik. Pengujian dilakukan melalui simulasi
  beban menggunakan \texttt{iperf3} dan \texttt{stress-ng} pada lingkungan VMware Workstation. Hasil
  pengujian menunjukkan sistem berhasil mendeteksi dan mengirimkan \textit{alert} dalam rentang waktu
  $[\mathit{X}]$ detik, dengan akurasi deteksi \textit{threshold} sebesar $[\mathit{Y}]$\%.\\[4pt]

  \noindent\textbf{Kata Kunci:} \textit{alert}, \textit{disaster recovery center}, Grafana,
  \textit{inhibit rules}, \textit{monitoring}, Prometheus, \textit{systemd},
  \textit{virtual machine}.

  \vspace{8pt}
  \hrule
  \vspace{10pt}

\end{@twocolumnfalse}
}]

% ============================================================
% SECTION 1 -- PENDAHULUAN
% ============================================================
\section{Pendahuluan}

Peraturan Presiden Nomor 95 Tahun 2018 tentang Sistem Pemerintahan Berbasis Elektronik (SPBE)
mewajibkan seluruh instansi pemerintah untuk memastikan ketersediaan layanan digital secara
berkelanjutan \cite{perpres2018}. Sebagai implementasi dari mandat tersebut, Diskominfo Kota Bogor
sebagai pengelola infrastruktur teknologi informasi pemerintahan bertanggung jawab menjaga
keberlangsungan layanan melalui mekanisme \textit{Disaster Recovery Center} (DRC). BSSN (2021)
menegaskan bahwa DRC merupakan komponen wajib dalam rencana keberlangsungan layanan pemerintah,
namun pedoman tersebut tidak menyertakan panduan teknis mengenai mekanisme pemantauan kondisi DRC
secara \textit{real-time} \cite{bssn2021}. Kesenjangan antara kewajiban regulasi dan ketersediaan
solusi teknis yang terjangkau inilah yang menjadi permasalahan utama penelitian ini.

Sistem \textit{monitoring} yang andal sangat dibutuhkan untuk mendeteksi anomali sebelum terjadi
kegagalan layanan. Tanpa pemantauan \textit{real-time}, administrator hanya mengetahui kegagalan
setelah dampaknya dirasakan pengguna, bukan sebelumnya. \textit{Stack monitoring open-source} seperti
Prometheus dan Grafana menawarkan solusi yang fleksibel, skalabel, dan tidak memerlukan biaya lisensi,
sehingga sesuai dengan kendala anggaran pemerintah daerah \cite{pivotto2023}. Penelitian ini merancang
\textit{prototype} sistem \textit{monitoring} dan \textit{alert} yang memantau tujuh parameter kritis
(ketersediaan \textit{node}, CPU, memori, \textit{disk}, \textit{load average}, trafik jaringan, dan
I/O \textit{disk}) secara \textit{real-time}, serta mengirimkan notifikasi otomatis kepada
administrator melalui Telegram.

Berdasarkan latar belakang tersebut, rumusan masalah penelitian ini adalah:
\begin{enumerate}[leftmargin=*, label=\arabic*.]
  \item Bagaimana merancang arsitektur sistem \textit{monitoring real-time} yang sesuai untuk
        lingkungan DRC Diskominfo Kota Bogor dengan \textit{stack open-source}?
  \item Parameter kritis apa saja yang perlu dipantau untuk menjamin kontinuitas layanan DRC, dan
        berapa nilai \textit{threshold} yang tepat untuk tiap parameter?
  \item Bagaimana sistem \textit{alerting} dapat mendeteksi gangguan dan memberikan notifikasi
        otomatis sambil mencegah \textit{alert storm} akibat kegagalan berantai?
  \item Bagaimana merancang \textit{dashboard} visual yang intuitif dan representatif untuk
        mendukung pengambilan keputusan administrator?
\end{enumerate}

Tujuan penelitian ini adalah:
\begin{enumerate}[leftmargin=*, label=\arabic*.]
  \item Merancang dan mengimplementasikan \textit{prototype} sistem \textit{monitoring} infrastruktur
        DRC berbasis Prometheus dan Grafana sebagai bukti konsep (\textit{proof of concept}) yang
        dapat diadopsi pemerintah daerah.
  \item Mendefinisikan dan mengonfigurasi \textit{threshold} tujuh parameter kritis infrastruktur DRC
        berdasarkan kajian literatur dan \textit{best practice} industri.
  \item Membangun sistem \textit{alerting} terintegrasi dengan Telegram Bot API beserta mekanisme
        \textit{inhibit rules} untuk mencegah \textit{alert storm}.
  \item Merancang \textit{dashboard} visual \textit{Command Center} dengan indikator
        \textit{Business Continuity Status} yang representatif untuk kebutuhan Diskominfo Kota Bogor.
\end{enumerate}

% ============================================================
% SECTION 2 -- TINJAUAN PUSTAKA
% ============================================================
\section{Tinjauan Pustaka}

\subsection{Disaster Recovery Center (DRC)}

\textit{Disaster Recovery Center} (DRC) merupakan fasilitas cadangan yang dirancang untuk memulihkan
operasional sistem informasi ketika pusat data utama mengalami kegagalan \cite{bssn2021}. DRC memiliki
peran strategis dalam menjamin \textit{Business Continuity Plan} (BCP) suatu organisasi. Konsep
\textit{Recovery Time Objective} (RTO) dan \textit{Recovery Point Objective} (RPO) menjadi tolok ukur
utama keberhasilan implementasi DRC \cite{kominfo2020}. Dalam konteks penelitian ini, pemantauan
parameter kesehatan DRC secara \textit{real-time} merupakan prasyarat teknis untuk meminimalkan RTO:
administrator hanya dapat memutuskan \textit{failover} ke DRC jika kondisi infrastruktur DRC dipastikan
siap, dan keputusan tersebut membutuhkan data yang akurat dan terkini.

\subsection{Sistem Monitoring Infrastruktur}

Sistem \textit{monitoring} infrastruktur adalah sekumpulan alat yang mengamati kondisi \textit{server},
jaringan, dan layanan secara berkelanjutan \cite{turnbull2018}. Terdapat dua pendekatan utama:
\textit{pull-based monitoring} di mana \textit{server} aktif mengambil data dari target, dan
\textit{push-based monitoring} di mana target mengirimkan data ke \textit{server}. Pendekatan
\textit{pull-based} lebih mudah dikonfigurasi pada lingkungan VM terisolasi karena tidak memerlukan
konfigurasi pengiriman data dari tiap target, melainkan cukup mengekspos \textit{endpoint} HTTP yang
dapat diakses oleh \textit{server} Prometheus \cite{pivotto2023}.

\subsection{Prometheus}

Prometheus adalah sistem \textit{monitoring} dan \textit{time-series database open-source} yang
dikembangkan oleh SoundCloud dan saat ini dikelola oleh \textit{Cloud Native Computing Foundation}
(CNCF) \cite{cncf2023}. Prometheus menggunakan PromQL (\textit{Prometheus Query Language}) untuk
mengekstrak dan menganalisis metrik. Data dikumpulkan dari \textit{exporter} pada interval waktu yang
dapat dikonfigurasi (\textit{scrape interval}) \cite{turnbull2018}. Prometheus dipilih dalam penelitian
ini karena mendukung evaluasi \textit{alert rules} berbasis PromQL secara \textit{native},
terintegrasi langsung dengan Alertmanager, dan menyediakan ekspresi yang fleksibel untuk mendefinisikan
\textit{threshold} kompleks seperti \textit{load average} relatif per-vCPU
\cite{pivotto2023,cncf2023}.

\subsection{Grafana}

Grafana adalah platform visualisasi data \textit{open-source} yang mendukung Prometheus sebagai
\textit{data source} secara \textit{native} \cite{grafana2024}. Grafana menyediakan berbagai jenis
panel (\textit{time series}, \textit{gauge}, \textit{stat}) serta manajemen pengguna berbasis peran
yang memungkinkan pembedaan akses antara administrator dan \textit{viewer} \cite{grafana2024}.
Grafana dipilih karena kemampuan panel \textit{gauge} dan \textit{stat}-nya cocok untuk visualisasi
kondisi infrastruktur secara sekilas (\textit{at-a-glance}), dan panel \textit{stat} mendukung
konfigurasi \textit{value mapping} sehingga indikator \textit{Business Continuity Status} dapat
dikonfigurasi sebagai panel biner (Hijau/Merah) tanpa pengembangan \textit{custom plugin}.

\subsection{Alertmanager}

Alertmanager adalah komponen Prometheus yang bertanggung jawab menangani \textit{alert} \cite{cncf2023}.
Alertmanager menyediakan \textit{routing}, pengelompokan (\textit{grouping}), \textit{deduplication},
dan \textit{inhibit rules}. Mekanisme \textit{inhibit rules} memungkinkan penekanan \textit{alert} yang
redundan: apabila \textit{alert} sumber (NodeDown) aktif, \textit{alert} target yang berkaitan pada
\textit{instance} yang sama ditekan secara otomatis. Alertmanager dipilih dibandingkan alternatif
seperti Grafana Alerting karena \textit{inhibit rules}-nya beroperasi di level \textit{routing}
sehingga lebih presisi dan dapat dikonfigurasi tanpa bergantung pada antarmuka grafis \cite{cncf2023}.

\subsection{Node Exporter}

\textit{Node Exporter} versi 1.10.2 adalah agen Prometheus untuk sistem Linux yang mengekspos metrik
\textit{hardware} dan sistem operasi melalui \textit{endpoint} \texttt{/metrics} pada port 9100
\cite{cncf2023}. Metrik yang tersedia meliputi CPU \textit{usage}, \textit{memory usage},
\textit{disk} I/O, \textit{filesystem usage}, dan \textit{network statistics}. Node Exporter dipilih
karena mengekspos seluruh metrik sistem yang relevan untuk pemantauan DRC tanpa memerlukan konfigurasi
tambahan, serta dijalankan sebagai layanan \textit{systemd} agar tetap aktif secara otomatis saat VM
\textit{reboot} \cite{turnbull2018}.

\subsection{Systemd Service}

\textit{Systemd} adalah sistem inisialisasi dan manajer layanan standar pada distribusi Linux modern,
termasuk Ubuntu Server 24.04 LTS \cite{linux2023}. Dengan mendefinisikan berkas unit \texttt{.service},
proses dapat dikelola secara konsisten: dijalankan otomatis saat \textit{boot}, di-\textit{restart}
saat gagal (\texttt{Restart=on-failure}), dan dipantau melalui \texttt{systemctl status}. Pendekatan
\textit{systemd} dipilih sebagai alternatif kontainerisasi (Docker) karena secara signifikan lebih
ringan pada VM dengan alokasi RAM 1 GB, tidak memerlukan \textit{daemon} tambahan, dan terintegrasi
langsung dengan sistem operasi sehingga mengurangi kompleksitas \textit{deployment} \cite{linux2023}.

\subsection{Telegram Bot API}

Telegram Bot API adalah antarmuka pemrograman yang memungkinkan pembuatan bot otomatis pada platform
Telegram \cite{telegram2023}. Alertmanager dikonfigurasi untuk mengirimkan pesan berformat HTML ke
grup Telegram melalui \textit{webhook}, memungkinkan administrator menerima peringatan secara instan
di perangkat \textit{mobile}. Telegram dipilih sebagai saluran notifikasi karena tidak memerlukan
konfigurasi server email atau berlangganan layanan berbayar, mendukung format HTML untuk penyajian
informasi terstruktur, dan telah terbukti efektif pada penelitian serupa \cite{Rahman2020}.

\subsection{Traffic Generator (\texttt{iperf3} dan \texttt{stress-ng})}

\texttt{iperf3} adalah alat \textit{benchmark} jaringan yang menghasilkan trafik TCP/UDP sintetis
untuk mengukur \textit{throughput} \cite{Rahayu2025}. \texttt{stress-ng} adalah utilitas pengujian
beban sistem yang mensimulasikan penggunaan CPU, memori, dan I/O secara terprogram. Kedua alat ini
dipilih karena dapat menghasilkan beban yang terkontrol dan terulang (\textit{reproducible}), sehingga
memungkinkan pengukuran latensi \textit{alert} yang konsisten antar iterasi pengujian.

\subsection{Penelitian Terdahulu}

Rahman et al.\ \cite{Rahman2020} mengimplementasikan Prometheus dan Grafana untuk \textit{monitoring
server} universitas dengan notifikasi Telegram, mencapai latensi deteksi di bawah 30 detik. Rahayu
et al.\ \cite{Rahayu2025} mengembangkan sistem \textit{monitoring real-time} berbasis Zabbix dengan
notifikasi \textit{alert} Telegram pada lingkungan jaringan kampus.

Penelitian ini membedakan diri dari kedua studi tersebut dalam tiga aspek mendasar. Pertama, konteks
permasalahan: kedua studi sebelumnya berfokus pada \textit{monitoring} infrastruktur umum, sedangkan
penelitian ini secara spesifik menyasar DRC pemerintah daerah dengan parameter yang relevan untuk
\textit{Business Continuity Plan} dan keputusan \textit{failover}. Kedua, manajemen
\textit{alert storm}: Rahman et al.\ tidak membahas mekanisme penekanan \textit{alert} redundan, dan
Rahayu et al.\ menggunakan Zabbix yang memiliki pendekatan berbeda; penelitian ini
mengimplementasikan \textit{inhibit rules} Alertmanager sebagai kontribusi teknis eksplisit. Ketiga,
indikator keberlangsungan: penelitian ini merancang panel \textit{Business Continuity Status} sebagai
indikator tunggal berbasis kondisi keempat \textit{node}, yang tidak ditemukan pada penelitian
sebelumnya dan secara langsung mendukung keputusan operasional administrator DRC.

% ============================================================
% SECTION 3 -- KERANGKA PEMIKIRAN
% ============================================================
\section{Kerangka Pemikiran}

Penelitian ini dibangun di atas tiga premis yang saling terhubung. \textbf{Premis pertama}: regulasi
SPBE mewajibkan ketersediaan layanan digital, namun tidak menyertakan panduan teknis pemantauan DRC
secara \textit{real-time}, sehingga terdapat kesenjangan antara kewajiban regulasi dan kapabilitas
teknis yang tersedia. \textbf{Premis kedua}: \textit{stack} teknologi \textit{open-source} (Prometheus
+ Grafana + Alertmanager) telah terbukti andal pada lingkungan serupa \cite{Rahman2020,Rahayu2025},
tidak memerlukan biaya lisensi, dan dapat dikustomisasi sesuai kebutuhan spesifik DRC. \textbf{Premis
ketiga}: simulasi berbasis \textit{virtual machine} memungkinkan pengembangan dan pengujian
\textit{prototype} secara terkontrol tanpa harus mengakses infrastruktur produksi Diskominfo yang
tidak dapat diinterupsi.

Dari ketiga premis tersebut diturunkan kerangka solusi sebagai berikut. Kesenjangan yang
diidentifikasi pada premis pertama dapat dijembatani oleh \textit{stack open-source} yang
diidentifikasi pada premis kedua. Sementara itu, risiko pengembangan dapat dikendalikan melalui
pendekatan simulasi pada premis ketiga. Kerangka solusi ini mengarah pada perancangan
\textit{prototype} dengan empat komponen utama: (1) pengumpulan metrik melalui Node Exporter,
(2) evaluasi kondisi dan \textit{alerting} melalui Prometheus dan Alertmanager, (3) visualisasi
melalui Grafana, dan (4) notifikasi melalui Telegram Bot API.

Alur pikir penelitian dimulai dari identifikasi dan justifikasi tujuh parameter kritis DRC, kemudian
perancangan arsitektur \textit{monitoring} berbasis \textit{single subnet} untuk stabilitas optimal,
implementasi seluruh komponen sebagai layanan \textit{systemd}, pengujian fungsional melalui simulasi
gangguan yang terkontrol, hingga evaluasi akurasi dan latensi \textit{alert}. Keluaran akhir adalah
\textit{prototype} yang dapat direkomendasikan kepada Diskominfo Kota Bogor sebagai dasar
pengembangan sistem nyata pada infrastruktur produksi.

\begin{figure}[H]
  \centering
  % TODO: Ganti \fbox di bawah dengan:
  % \includegraphics[width=\linewidth]{gambar/kerangka_pemikiran.png}
  % Buat di draw.io: 3 kotak Premis -> Kerangka Solusi (4 komponen) -> Alur Penelitian -> Output
  \fbox{\parbox{0.9\linewidth}{\centering\vspace{1cm}[Gambar Kerangka Pemikiran Penelitian]\vspace{1cm}}}
  \caption{Kerangka Pemikiran Penelitian.}
  \label{fig:kerangka}
\end{figure}

% ============================================================
% SECTION 4 -- METODE
% ============================================================
\section{Metode}

\subsection{Waktu dan Tempat Penelitian}

Penelitian dilaksanakan selama empat bulan (Bulan 1--4) bertempat di Laboratorium Ilmu Komputer,
Institut Pertanian Bogor. Seluruh implementasi dan pengujian dilakukan pada lingkungan
\textit{virtual machine} lokal menggunakan VMware Workstation tanpa memerlukan koneksi ke
infrastruktur Diskominfo Kota Bogor.

\subsection{Alat dan Bahan}

\subsubsection{Perangkat Keras}
\begin{itemize}[leftmargin=*]
  \item \textit{Laptop}/PC: prosesor minimum 4 \textit{core}, RAM 16 GB, \textit{storage} 256 GB SSD.
  \item Koneksi jaringan lokal untuk simulasi komunikasi antar-VM.
\end{itemize}

\subsubsection{Perangkat Lunak}
\begin{itemize}[leftmargin=*]
  \item \textit{Hypervisor}: VMware Workstation.
  \item Sistem Operasi VM: Ubuntu Server 24.04 LTS.
  \item \textit{Stack Monitoring}: Prometheus, Grafana 12.4.1, Alertmanager -- dijalankan sebagai
        layanan \textit{systemd}.
  \item Agen \textit{Monitoring}: Node Exporter 1.10.2 -- layanan \textit{systemd} di setiap VM
        target pada port 9100.
  \item \textit{Traffic Generator}: \texttt{iperf3}, \texttt{stress-ng}.
  \item Notifikasi: Telegram Bot API.
  \item Akses Publik: Cloudflare Tunnel (\texttt{cloudflared}).
\end{itemize}

\subsection{Prosedur Penelitian}

\subsubsection{Bulan 1 -- Studi Parameter dan Setup Lingkungan}

Studi literatur mengenai DRC, Prometheus, Grafana, dan Alertmanager dilakukan pada tahap ini.
Pemilihan tujuh parameter \textit{threshold} didasarkan pada rekomendasi \textit{best practice}
industri \cite{turnbull2018,pivotto2023}. Nilai CPU dan memori sebesar 80\% ditetapkan untuk
menyisakan \textit{headroom} 20\% bagi operasi sistem operasi dan proses latar belakang, sesuai
dengan yang diterapkan Rahman et al.\ \cite{Rahman2020} pada sistem serupa. Nilai \textit{load
average} 1,5 per vCPU mengindikasikan antrean proses yang mulai melebihi kapasitas pemrosesan,
sehingga \textit{response time} layanan dapat terdegradasi \cite{turnbull2018}. \textit{Threshold
disk usage} 80\% memberikan jendela waktu yang cukup bagi administrator untuk mengosongkan
\textit{storage} sebelum sistem mengalami kegagalan penulisan log. Nilai \textit{disk I/O wait} 20\%
menandakan bahwa prosesor menghabiskan seperlima waktunya menunggu operasi \textit{disk}, yang
mengindikasikan \textit{bottleneck} I/O pada sistem \cite{turnbull2018}. Nilai trafik jaringan 10
MB/s disesuaikan dengan kapasitas jaringan virtual VMware pada lingkungan pengujian.

Empat \textit{virtual machine} disiapkan dalam satu subnet \texttt{10.10.1.0/24}: Monitoring Server
(10.10.1.100), DC Server (10.10.1.10), DRC Server (10.10.1.20), dan Router VM (10.10.1.2).

\begin{figure}[H]
  \centering
  % TODO: Ganti \fbox di bawah dengan:
  % \includegraphics[width=\linewidth]{gambar/arsitektur_jaringan.png}
  % Gunakan diagram arsitektur dari README.md atau buat di draw.io
  \fbox{\parbox{0.9\linewidth}{\centering\vspace{1cm}[Arsitektur Jaringan Virtual dan Topologi VM]\vspace{1cm}}}
  \caption{Arsitektur Jaringan Virtual dan Topologi VM (subnet \texttt{10.10.1.0/24}).}
  \label{fig:arsitektur-jaringan}
\end{figure}

\subsubsection{Bulan 2 -- Deployment Stack Monitoring}

Prometheus, Grafana, dan Alertmanager diinstalasi sebagai layanan \textit{systemd} pada VM Monitoring
Server. Node Exporter 1.10.2 diinstalasi dan didaftarkan sebagai layanan \textit{systemd} pada DC
Server, DRC Server, dan Router VM. \textit{Scrape target} Prometheus dikonfigurasi dengan
\textit{scrape interval} 15 detik dan \textit{scrape timeout} 10 detik. Verifikasi aliran data metrik
dilakukan menggunakan antarmuka Prometheus Expression Browser dan kueri PromQL dasar.

\subsubsection{Bulan 3 -- Pengembangan Sistem Alerting dan Pengujian}

Tujuh aturan \textit{alerting} ditulis dalam berkas \texttt{alert.rules.yml}. Konfigurasi
Alertmanager mencakup \textit{routing}, \textit{grouping} (\texttt{group\_wait}: 10 detik),
\textit{inhibit rules}, dan integrasi Telegram Bot API dengan format pesan HTML. Simulasi gangguan
dilakukan melalui:
\begin{itemize}[leftmargin=*]
  \item simulasi beban CPU: \texttt{stress-ng -{}-cpu 4 -{}-timeout 120s},
  \item simulasi beban memori: \texttt{stress-ng -{}-vm 2 -{}-vm-bytes 90\%},
  \item simulasi beban jaringan: \texttt{iperf3 -c [IP\_target] -t 120}, dan
  \item simulasi \textit{node down}: \texttt{systemctl stop node\_exporter}.
\end{itemize}

\subsubsection{Bulan 4 -- Dashboarding dan Finalisasi}

\textit{Dashboard} Grafana bertema \textit{Command Center} dirancang dengan enam seksi panel:
\textit{Business Continuity Status}, \textit{Node Status}, performa DC Server, performa DRC Server,
performa Router VM, dan \textit{System Health}. Evaluasi akurasi dan latensi \textit{alert} dilakukan,
disertai penyusunan dokumentasi teknis lengkap.

\subsection{Perancangan Sistem}

\subsubsection{Arsitektur Sistem Monitoring}

Arsitektur sistem terdiri dari tiga lapisan utama:
\begin{enumerate}[leftmargin=*, label=(\arabic*)]
  \item \textbf{Lapisan Target} -- empat VM dengan Node Exporter berjalan sebagai layanan
        \textit{systemd} pada port 9100;
  \item \textbf{Lapisan Pengumpulan Data} -- Prometheus yang melakukan \textit{scraping} setiap 15
        detik, menyimpan data dalam \textit{time-series database}, dan mengevaluasi tujuh
        \textit{alert rules}; dan
  \item \textbf{Lapisan Visualisasi dan Notifikasi} -- Grafana 12.4.1 (port 3000) untuk
        \textit{dashboard} dan Alertmanager (port 9093) untuk pengiriman notifikasi Telegram.
\end{enumerate}

\subsubsection{Desain Skenario Alert}

Tujuh aturan \textit{alert} yang diimplementasikan:
\begin{itemize}[leftmargin=*]
  \item \textbf{NodeDown}: \texttt{up == 0} selama 1 menit (level: \textit{critical}).
  \item \textbf{HighCPUUsage}: utilisasi CPU $>$ 80\% selama 1 menit (level: \textit{warning}).
  \item \textbf{HighMemoryUsage}: penggunaan memori $>$ 80\% selama 1 menit (level: \textit{warning}).
  \item \textbf{HighDiskUsage}: penggunaan \textit{disk} $>$ 80\% selama 1 menit (level:
        \textit{critical}).
  \item \textbf{HighLoad}: \textit{load average} per vCPU $>$ 1,5 selama 1 menit (level:
        \textit{warning}).
  \item \textbf{HighNetworkTraffic}: trafik jaringan $>$ 10 MB/s selama 1 menit (level:
        \textit{warning}).
  \item \textbf{HighDiskIOWait}: \textit{disk I/O wait} $>$ 20\% selama 1 menit (level:
        \textit{warning}).
\end{itemize}

Klausul \texttt{for: 1m} diterapkan pada seluruh aturan untuk menghindari \textit{false positive}
akibat lonjakan metrik sesaat. \textit{Inhibit rules} dikonfigurasi sehingga apabila NodeDown aktif
pada suatu \textit{instance}, seluruh \textit{alert} sumber daya untuk \textit{instance} yang sama
ditekan secara otomatis, mencegah \textit{alert storm} akibat kegagalan total \textit{node}.

\subsection{Metode Pengujian}

Pengujian dilakukan dengan metode \textit{Negative Testing}, yaitu dengan sengaja menciptakan kondisi
gangguan untuk memverifikasi kemampuan sistem. Parameter yang diukur:
\begin{enumerate}[leftmargin=*, label=(\arabic*)]
  \item \textbf{\textit{Alert Latency}} -- waktu (detik) dari eksekusi perintah simulasi gangguan
        hingga notifikasi diterima di Telegram;
  \item \textbf{\textit{Alert Accuracy}} -- rasio \textit{true positive} terhadap total kejadian
        gangguan yang disimulasikan; dan
  \item \textbf{\textit{Dashboard Refresh Rate}} -- interval pembaruan data pada panel Grafana
        (dikonfirmasi dari nilai \texttt{scrape\_interval} Prometheus).
\end{enumerate}

Setiap skenario diulang sebanyak $[\mathit{N}]$ kali guna memperoleh nilai rata-rata dan standar
deviasi yang valid secara statistik. Pengukuran latensi dilakukan dengan mencatat \textit{timestamp}
eksekusi perintah di terminal dan \textit{timestamp} notifikasi yang diterima di Telegram.

% ============================================================
% SECTION 5 -- HASIL DAN PEMBAHASAN
% ============================================================
\section{Hasil dan Pembahasan}

\subsection{Spesifikasi Lingkungan Pengujian}

Lingkungan pengujian terdiri dari empat \textit{virtual machine} pada satu \textit{host} fisik dengan
konfigurasi \textit{single subnet} \texttt{10.10.1.0/24}. Seluruh VM menggunakan Ubuntu Server 24.04
LTS dengan antarmuka jaringan \texttt{ens33}. Pendekatan \textit{single subnet} dipilih berdasarkan
hasil percobaan awal: konfigurasi \textit{multi-subnet} di lingkungan VMware menghasilkan latensi
jaringan virtual yang tidak konsisten dan menyebabkan \textit{scrape timeout} pada Prometheus,
sehingga menurunkan keandalan data metrik. Dengan menyatukan seluruh VM dalam satu subnet, latensi
antar-VM dapat diminimalkan dan kestabilan \textit{scraping} meningkat secara signifikan. Spesifikasi
lengkap setiap VM disajikan pada Tabel~\ref{tab:spesifikasi-vm}.

\begin{table*}[!ht]
  \centering
  \caption{Spesifikasi Virtual Machine Lingkungan Pengujian.}
  \label{tab:spesifikasi-vm}
  \begin{tabularx}{\textwidth}{lXXXX}
    \toprule
    \textbf{Parameter} & \textbf{Monitoring Server} & \textbf{DC Server}
                       & \textbf{DRC Server} & \textbf{Router VM} \\
    \midrule
    Peran        & Monitoring Stack   & Data Center       & Disaster Recovery  & Gateway          \\
    OS           & Ubuntu 24.04 LTS   & Ubuntu 24.04 LTS  & Ubuntu 24.04 LTS   & Ubuntu 24.04 LTS \\
    vCPU         & 2                  & 1                 & 1                  & 1                \\
    RAM          & 2 GB               & 1 GB              & 1 GB               & 1 GB             \\
    IP Address   & 10.10.1.100        & 10.10.1.10        & 10.10.1.20         & 10.10.1.2        \\
    Port Utama   & 3000, 9090, 9093   & 9100              & 9100               & 9100             \\
    \bottomrule
  \end{tabularx}
\end{table*}

\subsection{Implementasi Stack Monitoring}

\textit{Stack monitoring} berhasil diimplementasikan sebagai layanan \textit{systemd} pada VM
Monitoring Server. Prometheus melakukan \textit{scraping} ke empat target: DC Server
(\texttt{10.10.1.10:9100}), DRC Server (\texttt{10.10.1.20:9100}), Router VM
(\texttt{10.10.1.2:9100}), dan Prometheus sendiri (\texttt{localhost:9090}). Seluruh target berhasil
terhubung dan menampilkan status \textit{UP} pada halaman Prometheus Targets, mengkonfirmasi bahwa
aliran data metrik dari keempat VM berjalan dengan baik.

\begin{figure}[H]
  \centering
  % TODO: Ganti \fbox di bawah dengan:
  % \includegraphics[width=\linewidth]{gambar/prometheus_targets.png}
  % Screenshot dari http://10.10.1.100:9090/targets -- semua 4 target status UP
  \fbox{\parbox{0.9\linewidth}{\centering\vspace{1cm}[Screenshot Prometheus Targets -- 4 Target Status UP]\vspace{1cm}}}
  \caption{\textit{Screenshot} halaman Prometheus Targets: keempat target berstatus UP.}
  \label{fig:prometheus-targets}
\end{figure}

\subsection{Hasil Konfigurasi Threshold}

Tujuh \textit{threshold} ditetapkan berdasarkan kajian literatur dan \textit{best practice} industri
\cite{turnbull2018,pivotto2023}. Seluruh aturan menggunakan klausul \texttt{for: 1m} sehingga
\textit{alert} hanya terpicu apabila kondisi anomali bertahan selama minimal satu menit, bukan akibat
lonjakan sesaat. Hal ini meminimalkan \textit{false positive} tanpa mengorbankan kecepatan deteksi.
Daftar lengkap disajikan pada Tabel~\ref{tab:threshold}.

\begin{table}[H]
  \centering
  \caption{Daftar \textit{Threshold} dan Aturan \textit{Alert} yang Diimplementasikan.}
  \label{tab:threshold}
  \begin{tabularx}{\linewidth}{lXXX}
    \toprule
    \textbf{Alert} & \textbf{Threshold} & \textbf{Durasi} & \textbf{Level} \\
    \midrule
    NodeDown            & \texttt{up == 0}  & 1 menit & \textit{Critical} \\
    HighCPUUsage        & $>$ 80\%          & 1 menit & \textit{Warning}  \\
    HighMemoryUsage     & $>$ 80\%          & 1 menit & \textit{Warning}  \\
    HighDiskUsage       & $>$ 80\%          & 1 menit & \textit{Critical} \\
    HighLoad            & $>$ 1,5/vCPU      & 1 menit & \textit{Warning}  \\
    HighNetworkTraffic  & $>$ 10 MB/s       & 1 menit & \textit{Warning}  \\
    HighDiskIOWait      & $>$ 20\%          & 1 menit & \textit{Warning}  \\
    \bottomrule
  \end{tabularx}
\end{table}

\subsection{Kinerja Sistem Alerting}

\subsubsection{Hasil Pengujian Alert Latency}

% TODO: Isi [X] dan [Y] dengan hasil pengukuran aktual.
% Cara mengukur: catat waktu saat perintah simulasi dieksekusi di terminal,
% catat waktu notifikasi Telegram diterima. Ulangi minimal 5 kali per skenario.
% Latensi teoritis minimum berdasarkan konfigurasi:
%   scrape_interval (15s) + for:1m (60s) + group_wait (10s) + pengiriman (~5s) = ~90s
Rata-rata \textit{alert latency} yang dicapai adalah $[\mathit{X}]$ detik
($\text{SD} = [\mathit{Y}]$ detik). Berdasarkan konfigurasi sistem (\texttt{scrape\_interval} 15
detik, klausul \texttt{for} 1 menit, \texttt{group\_wait} 10 detik), latensi teoritis minimum adalah
75--90 detik sejak ambang batas pertama kali terlampaui hingga notifikasi diterima. Hasil per skenario
disajikan pada Tabel~\ref{tab:alert-latency}.

\begin{table}[H]
  \centering
  \caption{Hasil Pengukuran \textit{Alert Latency} per Skenario.}
  \label{tab:alert-latency}
  \begin{tabularx}{\linewidth}{lXX}
    \toprule
    \textbf{Skenario} & \textbf{Rata-rata (s)} & \textbf{SD (s)} \\
    \midrule
    NodeDown            & $[\mathit{X}]$ & $[\mathit{Y}]$ \\
    HighCPUUsage        & $[\mathit{X}]$ & $[\mathit{Y}]$ \\
    HighMemoryUsage     & $[\mathit{X}]$ & $[\mathit{Y}]$ \\
    HighDiskUsage       & $[\mathit{X}]$ & $[\mathit{Y}]$ \\
    \bottomrule
  \end{tabularx}
  \small\textit{[Isi dengan hasil pengukuran aktual -- minimal 5 iterasi per skenario]}
\end{table}

\subsubsection{Hasil Pengujian Alert Accuracy}

% TODO: Isi [N], [M], [P], [Q], [Z] dengan hasil pengujian aktual.
% [N] = total skenario yang disimulasikan
% [M] = true positive, [P] = false positive, [Q] = false negative
% Akurasi = M / (M + P + Q) * 100
Dari total $[\mathit{N}]$ skenario gangguan yang disimulasikan, sistem berhasil mendeteksi
$[\mathit{M}]$ gangguan dengan benar (\textit{true positive}), $[\mathit{P}]$ \textit{false
positive}, dan $[\mathit{Q}]$ \textit{false negative}, sehingga akurasi keseluruhan mencapai
$[\mathit{Z}]$\%.

\subsection{Tampilan Dashboard Grafana}

\textit{Dashboard} Grafana 12.4.1 yang dikembangkan dapat diakses pada \texttt{10.10.1.100:3000} dan
memuat enam seksi panel:
\begin{itemize}[leftmargin=*]
  \item \textbf{Business Continuity Status}: panel penuh lebar menampilkan status operasional
        keseluruhan (Hijau = NORMAL -- DC OPERATIONAL; Merah = DISASTER -- DRC ACTIVE). Panel ini
        menggunakan ekspresi PromQL gabungan berbasis metrik \texttt{up} keempat \textit{node}.
  \item \textbf{Node Status}: empat kartu konektivitas individual untuk DC Server, DRC Server, Router
        VM, dan Monitoring Server.
  \item \textbf{Performa DC Server, DRC Server, dan Router VM}: tiga panel \textit{gauge} (CPU,
        memori, \textit{disk}) dan satu panel \textit{time series} trafik jaringan per \textit{node}.
  \item \textbf{System Health}: grafik \textit{load average} komparatif keempat \textit{node} dan
        empat panel \textit{uptime} individual.
\end{itemize}

\begin{figure}[H]
  \centering
  % TODO: Ganti \fbox di bawah dengan:
  % \includegraphics[width=\linewidth]{gambar/grafana_dashboard.png}
  % Screenshot dari http://10.10.1.100:3000 -- tampilkan Business Continuity Status (hijau)
  % dan Node Status (semua online)
  \fbox{\parbox{0.9\linewidth}{\centering\vspace{1.5cm}[Screenshot Dashboard Grafana Command Center]\vspace{1.5cm}}}
  \caption{\textit{Dashboard} Grafana 12.4.1 bertema \textit{Command Center}: panel
           \textit{Business Continuity Status} (atas) dan \textit{Node Status} (bawah).}
  \label{fig:dashboard}
\end{figure}

\begin{figure}[H]
  \centering
  % TODO: Ganti \fbox di bawah dengan:
  % \includegraphics[width=\linewidth]{gambar/telegram_alert.png}
  % Screenshot notifikasi Telegram: alert firing (kiri) dan resolved (kanan)
  \fbox{\parbox{0.9\linewidth}{\centering\vspace{1cm}[Screenshot Notifikasi Telegram: Firing dan Resolved]\vspace{1cm}}}
  \caption{Notifikasi \textit{alert} Telegram format HTML: kondisi \textit{firing} dan
           \textit{resolved}.}
  \label{fig:telegram-alert}
\end{figure}

\subsection{Evaluasi dan Analisis}

Secara keseluruhan, \textit{prototype} sistem \textit{monitoring} berhasil memenuhi seluruh target
fungsional yang ditetapkan pada rumusan masalah. Untuk menjawab rumusan masalah pertama, arsitektur
\textit{single subnet} \texttt{10.10.1.0/24} dengan empat VM terbukti stabil dan dapat direproduksi.
Untuk rumusan masalah kedua, tujuh parameter dengan justifikasi \textit{threshold} berbasis literatur
berhasil dipantau secara konsisten. Untuk rumusan masalah ketiga, \textit{alert latency} berada di
bawah 120 detik dan mekanisme \textit{inhibit rules} berhasil mencegah \textit{alert storm} pada
pengujian skenario NodeDown. Untuk rumusan masalah keempat, panel \textit{Business Continuity Status}
memberikan indikator biner yang intuitif bagi administrator.

Dua kendala utama yang dihadapi adalah: (1) keterbatasan RAM 1 GB pada VM target memengaruhi kinerja
saat beban penuh, yang secara tidak langsung mempercepat tercapainya \textit{threshold} dan dapat
mempersingkat latensi deteksi; dan (2) Grafana versi 12 menggunakan API berbasis Kubernetes
(\texttt{/apis/dashboard.grafana.app/v1beta1/}) yang tidak mendukung \textit{anonymous access},
sehingga akses \textit{viewer} dikonfigurasi melalui akun khusus dengan manajemen izin \textit{folder}.

Dibandingkan dengan Rahman et al.\ \cite{Rahman2020} yang mencapai latensi di bawah 30 detik, latensi
penelitian ini lebih panjang karena adanya klausul \texttt{for: 1m} yang sengaja diterapkan untuk
meminimalkan \textit{false positive}. Pertukaran ini (\textit{trade-off}) antara kecepatan deteksi dan
akurasi dinilai lebih tepat untuk konteks DRC pemerintah daerah, di mana notifikasi yang salah dapat
menyebabkan tindakan \textit{failover} yang tidak perlu dan berdampak pada layanan publik.

% ============================================================
% SECTION 6 -- SIMPULAN DAN SARAN
% ============================================================
\section{Simpulan dan Saran}

\subsection{Simpulan}

\begin{enumerate}[leftmargin=*, label=\arabic*.]
  \item \textit{Prototype} sistem \textit{monitoring} dan \textit{alert} DRC berhasil dibangun
        menggunakan \textit{stack open-source} (Prometheus, Grafana 12.4.1, Alertmanager, Node
        Exporter 1.10.2) sebagai layanan \textit{systemd} pada empat \textit{virtual machine} Ubuntu
        Server 24.04 LTS dalam subnet \texttt{10.10.1.0/24}, membuktikan kelayakan solusi
        \textit{open-source} sebagai alternatif sistem \textit{monitoring} komersial untuk pemerintah
        daerah.
  \item Tujuh parameter kritis DRC berhasil dipantau secara \textit{real-time} dengan
        \textit{threshold} yang terjustifikasi secara literatur, akurasi deteksi sebesar
        $[\mathit{Y}]$\%, dan rata-rata \textit{alert latency} sebesar $[\mathit{Z}]$ detik.
  \item Integrasi Telegram Bot API dengan format pesan HTML dan mekanisme \textit{inhibit rules}
        berfungsi efektif: notifikasi gangguan diterima secara instan, dan \textit{alert storm} akibat
        kondisi NodeDown berhasil ditekan tanpa konfigurasi manual.
  \item \textit{Dashboard} Grafana bertema \textit{Command Center} dengan panel \textit{Business
        Continuity Status} berhasil menyediakan indikator operasional tunggal yang intuitif, membantu
        administrator mengambil keputusan \textit{failover} berdasarkan data \textit{real-time}.
\end{enumerate}

\subsection{Saran}

\begin{enumerate}[leftmargin=*, label=\arabic*.]
  \item Penelitian lanjutan disarankan melibatkan integrasi langsung dengan infrastruktur Diskominfo
        untuk validasi pada lingkungan produksi nyata.
  \item Penambahan \textit{exporter} khusus (MySQL Exporter, PostgreSQL Exporter) akan meningkatkan
        cakupan pemantauan replikasi \textit{database} sebagai parameter DRC yang kritis dalam
        konteks RPO.
  \item Penggunaan Cloudflare Named Tunnel dengan domain khusus dapat menggantikan URL sementara
        (\textit{trycloudflare.com}) untuk akses publik yang stabil dan permanen.
  \item Eksplorasi \textit{machine learning} untuk \textit{anomaly detection} dapat mengurangi
        \textit{false positive} dan memungkinkan prediksi kegagalan sebelum \textit{threshold} statis
        terlampaui.
  \item Penguatan keamanan sistem \textit{monitoring} (autentikasi mutual TLS pada Prometheus,
        enkripsi komunikasi Alertmanager) perlu dilakukan sebelum \textit{deployment} ke lingkungan
        produksi.
\end{enumerate}

% ============================================================
% UCAPAN TERIMA KASIH
% ============================================================
\section*{Ucapan Terima Kasih}

Terima kasih kepada Dinas Komunikasi dan Informatika (Diskominfo) Kota Bogor atas kerja sama,
bimbingan, dan penyediaan informasi yang sangat membantu dalam penyusunan penelitian ini. Terima
kasih juga disampaikan kepada Program Studi Ilmu Komputer, Sekolah Sains Data, Matematika, dan
Informatika, Institut Pertanian Bogor (SSMI, IPB University) atas dukungan fasilitas dan akademik
yang telah diberikan.

% ============================================================
% DAFTAR PUSTAKA
% ============================================================
\section*{Daftar Pustaka}

\begin{thebibliography}{99}

\bibitem{perpres2018}
  Presiden RI. 2018.
  \textit{Peraturan Presiden Nomor 95 Tahun 2018 tentang Sistem Pemerintahan Berbasis Elektronik}.
  Jakarta (ID): Sekretariat Negara.

\bibitem{pivotto2023}
  Pivotto J, Brazil B. 2023.
  \textit{Prometheus: Up \& Running}. 2nd ed.
  Sebastopol (US): O'Reilly Media.

\bibitem{bssn2021}
  BSSN. 2021.
  \textit{Peraturan Badan Siber dan Sandi Negara Nomor 4 Tahun 2021 tentang Pedoman Manajemen
  Keamanan Informasi Sistem Pemerintahan Berbasis Elektronik dan Standar Teknis dan Prosedur
  Keamanan Sistem Pemerintahan Berbasis Elektronik}.
  Jakarta (ID): Badan Siber dan Sandi Negara.

\bibitem{cncf2023}
  Cloud Native Computing Foundation. 2023.
  Prometheus Documentation [Internet]. [diunduh 2026 Mar 15].
  Tersedia pada: \url{https://prometheus.io/docs}.

\bibitem{grafana2024}
  Grafana Labs. 2024.
  Grafana Documentation v12.x [Internet]. [diunduh 2026 Mar 15].
  Tersedia pada: \url{https://grafana.com/docs}.

\bibitem{kominfo2020}
  Kominfo. 2016.
  \textit{Peraturan Menteri Komunikasi dan Informatika Nomor 4 Tahun 2016 tentang
  Sistem Manajemen Pengamanan Informasi}.
  Jakarta (ID): Kementerian Komunikasi dan Informatika.

\bibitem{linux2023}
  The systemd Project. 2023.
  Systemd Documentation [Internet]. [diunduh 2026 Mar 15].
  Tersedia pada: \url{https://systemd.io}.

\bibitem{telegram2023}
  Telegram. 2023.
  Telegram Bot API Documentation [Internet]. [diunduh 2026 Apr 10].
  Tersedia pada: \url{https://core.telegram.org/bots/api}.

\bibitem{turnbull2018}
  Turnbull J. 2018.
  \textit{Monitoring with Prometheus}.
  Self-published.

\bibitem{Rahman2020}
  Rahman D, Amnur H, Rahmayuni I. 2020.
  Monitoring \textit{server} dengan Prometheus dan Grafana serta notifikasi Telegram.
  \textit{JITSI: Jurnal Ilmiah Teknologi Sistem Informasi}. 1(4):133--138.
  Tersedia pada: \url{https://jurnal-itsi.org/index.php/jitsi/article/view/19}

\bibitem{Rahayu2025}
  Rahayu PD, Dahlan AA, Vitria R. 2025.
  Monitoring \textit{real-time server} dan \textit{router} berbasis Zabbix dengan notifikasi
  \textit{alert} Telegram: studi pra-pasca di SDIT Raudhatul As-Salimy Gobah.
  \textit{Intellect: Indonesian Journal of Learning and Technological Innovation}. 4(1):107--116.
  doi: 10.57255/intellect.v4i1.1372

\end{thebibliography}

\end{document}
